\section{Introduction}
\label{sec:intro}

A marketplace has to do four things: let sellers publish what they offer, let buyers find it, carry an order from buyer to seller, and settle payment in a way both sides can accept. Hosted marketplaces do all four in one place. That place can see every order, hold every payment, and remove any seller. OpenBazaar's premise is that these functions can be split so that no single operator is needed for any of them, and so that the parties who must learn something (the seller who ships, the moderator who rules on a dispute) learn it only when they need to.

\subsection{The original OpenBazaar}

OpenBazaar began in 2014 as a prototype peer-to-peer market and was developed as \code{openbazaar-go} and a desktop client until late 2020~\citep{openbazaargo}. A node published its store as an IPFS directory addressed by IPNS~\citep{ipfs}. Listings were signed by the node's identity key. Messages traveled over libp2p when both peers were online; when the recipient was offline, the sender stored the message on its own node or an external storage host and left a pointer in the distributed hash table. Search was a separate hosted index that nodes notified on publish. The wallet supported several cryptocurrencies. Moderated purchases used a 2-of-3 multisignature script over buyer, vendor, and moderator keys, with a default timeout of 1080 hours.

That design had strengths the redesign keeps: listings were signed and an order committed to a hash of the listing it bought, so later edits could not rewrite a deal; the moderator could not take funds alone. It also had properties the redesign changes:
\begin{itemize}
  \item A store was available only while its own node, or a pinning service, was online.
  \item The libp2p peer identity was long-lived and observable, and peers dialed each other directly, so the network layer linked a store identity to IP addresses.
  \item The default search index was a single hosted service. Delisting from it removed a store from practical discovery.
  \item Spending a 2-of-3 P2SH or P2WSH output reveals the full script, including the moderator's public key, so every moderated purchase was recognizable on chain and attributable to its moderator.
  \item A phone could not run the daemon in the background, so a vendor without an always-on desktop missed orders.
  \item The multi-coin wallet made wallet maintenance the dominant engineering cost.
\end{itemize}

\subsection{This paper}

The redesign keeps the contract between buyer, vendor, and moderator and replaces the machinery underneath. Its main ideas are:
\begin{enumerate}
  \item \textbf{A public catalog, private orders.} Store and listing records are public, signed, and content-addressed, and replicated by any node (\S\ref{sec:catalog}). Orders are not catalog records; each is a letter sealed to one vendor and carried to that vendor's own mailbox (\S\ref{sec:orders}). Vendors label their own records, and buyers filter what they see on their own devices (\S\ref{sec:labels}).
  \item \textbf{Clients that only talk to their own node.} The user-facing app sends marketplace requests to its owner's node, which verifies records and delivers letters (\S\ref{sec:architecture}). A phone reaches that node as in \S\ref{sec:remote}.
  \item \textbf{Direct, per-seller payments.} A buyer pays each vendor in a separate transaction, never batching two sellers together (\S\ref{sec:payments}).
  \item \textbf{Scriptless two-party escrow with a bounded moderator.} The escrow output is an ordinary-looking Taproot key-path output. Every non-cooperative outcome is presigned by both parties before the buyer funds it. The moderator only reveals one secret scalar that completes one of two presigned dispute outcomes (\S\ref{sec:escrow}).
\end{enumerate}

The contributions are: (i) a threat model for a marketplace with this split (\S\ref{sec:threats}); (ii) a precise description of the two-party escrow protocol, its message rounds, the checks each party and the moderator must run, and its block-height timeline (\S\ref{sec:escrow}); (iii) argued security properties, each with its assumptions, and the attacks the checks exist to stop (\S\ref{sec:security}); (iv) a table of what every party and observer learns (\S\ref{sec:privacy}); and (v) measurements of a reference implementation (\S\ref{sec:evaluation}) and a statement of what the design does not achieve (\S\ref{sec:limitations}).

A reference implementation in Rust, with a web and desktop client and a native phone app for iOS and Android, is available at \url{https://github.com/OpenBazaar/openbazaar-rs}. This paper describes what that implementation does.
